\chapter{Les différents types de systèmes de fichiers}
	\section{Systèmes de fichiers locaux}
\begin{description}
    \item[Minix] Premier système de fichiers utilisé sous Linux.
    \item[ext2] Système de fichiers standard de Linux.
    \item[msdos] Système FAT16 utilisé sous MS-DOS et Windows.
    \item[vfat] Système FAT32 utilisé sous Windows.
    \item[ntfs] Système de fichiers de Windows NT.
\end{description}

	\section{Systèmes de fichiers réseau}
SMB

Système de fichiers réseau utilisant le protocole SMB de Microsoft.

NFS

Système de fichiers réseau développé par Sun Microsystems.

	\section{Système de fichiers optique}
ISO9660

Système de fichiers utilisé sur les CD-ROM.

	\section{Les systèmes de fichiers journalisés}
Les systèmes de fichiers journalisés assurent la cohérence des métadonnées grâce à un mécanisme de journalisation.

Exemples :
\begin{itemize}[label=$\bullet$]
    \item ext3
    \item ReiserFS
    \item XFS
\end{itemize}
En cas de panne ou de redémarrage après un incident, le système relit le journal afin de restaurer un état cohérent du système de fichiers.